\documentclass[11pt,a4paper]{article}

\usepackage[utf8]{inputenc}
\usepackage[T1]{fontenc}
\usepackage[brazil]{babel}
\usepackage{lmodern}
\usepackage[margin=2.5cm]{geometry}
\usepackage{amsmath,amssymb}
\usepackage{graphicx}
\usepackage{booktabs}
\usepackage{enumitem}
\usepackage{xcolor}
\usepackage{listings}
\usepackage[hidelinks]{hyperref}

\lstset{
  language=Python,
  basicstyle=\ttfamily\small,
  keywordstyle=\color{blue!60!black},
  commentstyle=\color{gray},
  stringstyle=\color{orange!70!black},
  columns=fullflexible,
  frame=single,
  rulecolor=\color{gray!40},
  aboveskip=0.8em,
  belowskip=0.8em,
}

\newcommand{\code}[1]{\texttt{#1}}

\title{Exercício 1: Tópicos Especiais em Aprendizado por Reforço}
\author{COS767 - 2026/3}
\date{}

\begin{document}
\maketitle

\section{Cenário}

Vamos programar um robô de limpeza que divide o espaço com um cachorro territorialista! O cachorro
se movimenta pela sala ``marcando território'' e cabe ao robô limpar a bagunça. O robô é penalizado
sempre que a sala está suja, então ele deve passar onde o cachorro molhou o mais rápido possível!
Para dificultar a história, o cachorro também é brincalhão, e se percebe o robô se movendo, ele o
persegue, e pula em cima dele. Isso faz com que o robô precise reiniciar, ficando $k$ rodadas
parado enquanto o cachorro anda livremente pela sala.

\begin{figure}[h]
  \centering
  \begin{minipage}{0.49\textwidth}
    \centering
    \includegraphics[width=\textwidth]{figuras/sala_normal.png}\\[2pt]
    (a)
  \end{minipage}\hfill
  \begin{minipage}{0.49\textwidth}
    \centering
    \includegraphics[width=\textwidth]{figuras/sala_brincadeira.png}\\[2pt]
    (b)
  \end{minipage}
  \caption{A sala, contendo uma mesa no centro. (a) Cachorro em modo normal, olhando para a
    direita: as células em rosa indicam o seu campo de visão. (b) Cachorro em modo brincadeira
    (``!''), perseguindo o robô.}
  \label{fig:sala}
\end{figure}

\section{Descrição do MDP}

\paragraph{Sala.} Uma grade de $9 \times 5$ células com uma mesa de $3 \times 1$ no centro, que
bloqueia o movimento e a visão. Ou seja, há 42 células livres para o robô e para o cachorro.

\paragraph{Estado.} Posição do robô, posição do cachorro, orientação do cachorro ($N$, $L$, $S$ ou
$O$), modo do cachorro (normal ou brincadeira) e posição do xixi (ou sala limpa). O cachorro, sendo
territorialista, só faz xixi novamente quando nota que a sala está limpa, ou seja, há no máximo um
xixi na sala por vez. Obtemos então $2 \cdot 42 \cdot 42 \cdot 4 \cdot 43 = 606.816$ estados.

\paragraph{Ações.} Movimentar o robô uma célula para $N$, $L$, $S$ ou $O$, ou ficar parado. Um
movimento contra a parede ou a mesa não surte efeito.

\paragraph{Recompensa.} $r(s) = -1$ sempre que há xixi na sala e $r(s) = 0$ caso contrário.

\paragraph{Uma rodada.} A partir do estado $s$ e da ação $a$, os eventos acontecem na seguinte
ordem:
\begin{enumerate}
  \item O robô executa $a$, uma ação de movimento. Para limpar o xixi, basta ele terminar a rodada
    na mesma célula.
  \item Se o cachorro está em modo normal, ele percebe o robô quando o robô se moveu nesta rodada e
    (i) a nova posição do robô está no campo de visão do cachorro, ou (ii) o robô entrou na célula
    do cachorro. O campo de visão é um cone de $90^\circ$ na direção em que o cachorro olha,
    limitado pelo obstáculo central, a mesa. Ao perceber o robô, o cachorro entra no modo
    brincadeira já nesta rodada.
  \item O cachorro se move.
    \begin{itemize}
      \item \textbf{Modo normal.} Com probabilidade $1-p$, anda uma célula na direção que está
        virado. Se está de frente para uma barreira, o cachorro dá meia-volta. Com probabilidade
        $p$, o cachorro para e vira para uma das outras três direções, escolhida uniformemente, e,
        se a sala está limpa, faz xixi na própria célula.
      \item \textbf{Modo brincadeira.} Se o cachorro tem linha de visão até o robô (em qualquer
        direção), avança até duas células em direção a ele, incluindo diagonais, e passa a olhar
        para ele. Caso contrário, corre até duas células na direção em que estava olhando, ou seja,
        na direção em que viu o robô pela última vez. Se ele chega a um obstáculo sem ver o robô,
        volta ao modo normal.
    \end{itemize}
  \item Se o cachorro alcança a célula do robô, há uma captura. O cachorro volta ao modo normal e o
    robô é desligado! Enquanto o robô reinicia, o cachorro anda $k$ rodadas em modo normal (e pode
    fazer xixi).
\end{enumerate}

As $k$ rodadas de uma captura formam uma única transição, com recompensa $r(s)$. Dessa forma, a
propriedade de Markov é mantida: cada valor de $k$ apenas define uma distribuição para o próximo
estado, e o robô não precisa conhecer o caminho que o cachorro fez.

\paragraph{Parâmetros.} $p = 0{,}2$. Os experimentos devem usar pelo menos
$k \in \{0, 1, 2, 4, 6, 10\}$ e dois fatores de desconto: $\gamma = 0{,}7$ (curto prazo) e
$\gamma = 0{,}99$ (longo prazo).

\section{Código fornecido}

O repositório já contém o ambiente, o modelo e a animação. O projeto é em Python, e usa o
gerenciador de pacotes \code{uv}. A instalação pode ser feita com \code{uv sync}. Se o(a)
estudante tiver a preferência de outro gerenciador de pacotes ou linguagem de programação, é
permitido o uso de ferramentas de IA para a conversão do código. Os componentes a serem
implementados como parte do exercício, porém, devem ser feitos pela própria pessoa. Por isso a
maior parte do código auxiliar já está disponibilizada, para que só se preocupe com os experimentos
e algoritmos relevantes à disciplina.

Para ajudar a depurar as implementações, use a função \code{animate(policy, env, steps, filename)}
para salvar um GIF da sua política em ação (badum tss)!

\begin{center}
\begin{tabular}{@{}ll@{}}
  \toprule
  Módulo & Conteúdo \\
  \midrule
  \code{cleaning\_robot/env.py} & \code{CleaningEnv}: dinâmica do cachorro e interface do MDP \\
  \code{cleaning\_robot/model.py} & \code{all\_states}, \code{state\_index}, \code{TabularModel} \\
  \code{cleaning\_robot/policy.py} & \code{PolicyMixin}, \code{GreedyPolicy}, \code{RandomPolicy} \\
  \code{cleaning\_robot/animate.py} & \code{animate} (GIF de uma política) e \code{rollout} (simulação) \\
  \code{cleaning\_robot/dp.py} & \code{TabularPolicy}, \code{tabulate} e os esqueletos dos itens~1 e~2 \\
  \code{cleaning\_robot/qlearning.py} & esqueleto do item~3 \\
  \bottomrule
\end{tabular}
\end{center}

\begin{lstlisting}
from cleaning_robot.env import CleaningEnv
from cleaning_robot.model import TabularModel, state_index
from cleaning_robot.policy import GreedyPolicy
from cleaning_robot.animate import animate, rollout

env = CleaningEnv(k=2)                  # p = 0.2
s = env.random_init(seed=0)             # estado inicial aleatorio
env.reward(s)                           # r(s)
env.next_state_probs(s, a)              # {s': P(s'|s,a)}
env.next_state_sample(s, a, seed)       # s' ~ P(.|s,a)

model = TabularModel(env)               # model.P[a][i, j] = P(j|i,a), model.R[i] = r(i)
i = state_index(env, s)                 # linha de s em P e R

animate(GreedyPolicy(env), env, steps=80, filename="greedy.gif")
\end{lstlisting}

\paragraph{Modelo.} Para facilitar uma implementação matricial dos algoritmos model-based,
fornecemos a classe \code{TabularModel}, que enumera os estados com \code{state\_index} e guarda o
MDP como uma matriz esparsa $n \times n$ por ação, \code{model.P[a]}, e o vetor de recompensas
\code{model.R}. Pode valer a pena salvar as matrizes em disco (\code{scipy.sparse.save\_npz}).

\paragraph{Políticas.} Toda política deve herdar \code{PolicyMixin} e implementar
\code{act(state, seed=None)}, que devolve uma ação. \code{TabularPolicy} guarda uma ação por índice
de estado; \code{tabulate(model, policy)} produz essa tabela para qualquer política
determinística. A política gulosa (\code{GreedyPolicy}) segue o caminho mais curto até o xixi e
ignora o cachorro; a política aleatória (\code{RandomPolicy}) escolhe uma ação uniforme.

\section{Funções e classes a implementar}
\label{sec:fc}

\begin{itemize}
  \item \code{policy\_evaluation(model, policy, gamma)} em \code{dp.py}: o valor $V^\pi$ de uma
    política, usando o modelo. A política é uma ação por estado (determinística) ou uma matriz de
    probabilidades $n \times 5$.
  \item \code{ValueIteration(model, gamma)} e \code{PolicyIteration(model, gamma)} em
    \code{dp.py}.
  \item \code{QLearning(env, gamma, n\_steps, seed)} em \code{qlearning.py}. Lembre que o
    Q-learning é model-free e só pode usar \code{random\_init}, \code{next\_state\_sample} e
    \code{reward}, não \code{model.P}.
\end{itemize}

As classes já têm um esqueleto, e herdam de \code{PolicyMixin}, mas é permitido mexer no código já
feito, inclusive nos construtores para incluir parâmetros adicionais.

\section{Tarefas}

Para os itens que seguem abaixo, você deve definir e justificar suas métricas de desempenho. Vamos
querer medir convergência para o ótimo, comparar políticas, etc.

Os itens 1 e 2 devem ser feitos para $\gamma = 0{,}7$ e $\gamma = 0{,}99$.

\begin{enumerate}
  \item \textbf{Avaliação de políticas.} Implemente \code{policy\_evaluation} e avalie a política
    gulosa e a política aleatória para cada $k$.
    \begin{enumerate}[label=(\alph*)]
      \item Quanto uma política perde em relação à outra, como medir isso?
      \item Como o valor de cada uma varia com $k$? E com $\gamma$?
      \item Qual abordagem é mais apropriada neste cenário: resolver o sistema diretamente ou por
        um método iterativo? Justifique comparando tempo e memória.
    \end{enumerate}

  \item \textbf{Policy iteration e value iteration.} Implemente os dois métodos e faça gráficos do
    tempo de convergência e do número de iterações em função de $k$.
    \begin{enumerate}[label=(\alph*)]
      \item O tempo de convergência muda muito com $k$? E o número de iterações? Por quê?
      \item A política ótima muda com $k$? Como medir isso?
      \item Compare a política ótima com as políticas do item 1.
    \end{enumerate}

  \item \textbf{Q-learning.} Escolha os dois cenários $(k, \gamma)$ mais interessantes e execute o
    Q-learning.
    \begin{enumerate}[label=(\alph*)]
      \item Como inicializar Q? Considere que só existem recompensas não-positivas. Isso tem alguma
        consequência em inicializar tudo como zero? Dica: a inicialização pode usar o valor de
        $\gamma$.
      \item Compare com os métodos baseados em modelo: qualidade da política aprendida ao longo do
        treino, número de passos de interação e tempo de execução.
    \end{enumerate}
\end{enumerate}

\section{Entrega}

Código e um relatório em PDF com os gráficos e as respostas às perguntas. Por favor, use o grupo do
WhatsApp para tirar dúvidas ou relatar problemas com o código fornecido. É a primeira vez que
elaboramos esse exercício.

\section{Uso de IA}

A formulação do problema e a especificação das interfaces foram feitas sem IA. O código, porém, foi
em sua maioria gerado automaticamente por uma LLM + harness. Como comentado anteriormente, se você
quiser converter o código fornecido para outra linguagem da sua escolha, sinta-se livre para
fazê-lo também com assistência de IA. Usar IA para tirar dúvidas do código também é permitido. No
entanto, não use IA para nenhuma outra coisa, em particular, na implementação do que está listado na
Seção~\ref{sec:fc} e na lógica dos experimentos.

\end{document}
